% MT050 / explanation — canonical editable content.
预测值\(f(x)\)用于准备单车，真实借车量为\(y\)。先固定残差方向\(r=f(x)-y\)：\(r<0\)表示准备少了、会缺车；\(r>0\)表示多备了车。公司更怕缺车，因此\textbf{相同误差大小下，负残差应付出更高代价}。


可设计非对称平方损失。为避免与线性模型的偏置\(b\)混淆，把两侧损失系数称为\(c_-\)和\(c_+\)：
\[
L(r)=\begin{cases}c_-r^2,&r<0,\\c_+r^2,&r\ge0,\end{cases}
\qquad c_->c_+>0.
\]


取教学参数\(c_-=4,c_+=1\)。少准备1辆，\(r=-1\)，损失为4；多准备1辆，\(r=1\)，损失为1。少2辆损失16，多2辆损失4。因此模型宁可容忍一些多备车，也会更努力避免同等大小的缺车。


画图时横轴标\(r=\text{预测}-\text{真实}\)，纵轴标\(L(r)\)。两半抛物线都在\((0,0)\)接上，左半边更陡，右半边较缓；\((-1,4)\)与\((1,1)\)是方便定位的点。两侧在0处斜率均为0，接点连续且光滑。


\StudyImage{CS5489/exam-review/Midterm/assets/answer-asymmetric-loss.png}


最后把每条记录的损失加起来训练：\(\min_{w,b}\sum_iL(w^Tx_i+b-y_i)\)。不对称代价改变了最优预测的取舍，但不保证每天绝不缺车。也可以设计非对称绝对损失，只要残差方向、贵的一侧及图形一致。


\textbf{出处：} 2023A期中 Q12，原卷第9页（10分）；2025A期中 Q12，原卷第9页（10分）


\textbf{答案依据：} 2023A期中Q12随卷方案；2025A Q12复用情境与要求。图中原系数a=1、b=4对应解析中的\(c_+=1,c_-=4\)，是教学作图参数。
